
\documentclass[12pt,a4paper]{article}

% Packages
\usepackage[T1]{fontenc}
\usepackage[utf8]{inputenc}
\usepackage[polish]{babel}
\usepackage{lmodern}
\usepackage{geometry}
\geometry{margin=2.3cm}
\usepackage{amsmath,amssymb,siunitx}
\sisetup{locale=DE, per-mode=symbol}
\usepackage{circuitikz}
\usepackage{tikz}
\usepackage{enumitem}
\usepackage{booktabs}
\usepackage{microtype}
\usepackage{hyperref}

\hypersetup{
    colorlinks=true,
    linkcolor=blue,
    urlcolor=blue
}

\setlength{\parskip}{0.45em}
\setlength{\parindent}{0pt}
\allowdisplaybreaks

% Custom commands
\newcommand{\jw}{\mathrm{j}\omega}
\newcommand{\Z}{\underline{Z}}
\newcommand{\U}{\underline{U}}
\newcommand{\I}{\underline{I}}
\newcommand{\Htr}{\underline{H}}

\title{\textbf{Praca domowa}\\Obwody prądu sinusoidalnego i filtry}
\author{}
\date{}

\begin{document}

\maketitle
\tableofcontents
\newpage

% ==================================================
\section{Wykresy wskazowe rezystora i kondensatora}
% ==================================================

\subsection*{Treść zadania}

Narysować wykres wskazowy dla rezystora $R$ i kondensatora $C$ połączonych:
\begin{enumerate}[label=\alph*)]
    \item szeregowo,
    \item równolegle.
\end{enumerate}

\subsection*{Najważniejsze wzory i teoria}

Dla przebiegów sinusoidalnych przyjmujemy zapis zespolony.
Impedancje elementów wynoszą:

\[
\Z_R=R,\qquad
\Z_C=\frac{1}{\jw C}
=-\mathrm{j}\frac{1}{\omega C}.
\]

W rezystorze napięcie i prąd są w fazie.
W kondensatorze prąd wyprzedza napięcie o $90^\circ$.

Wykres wskazowy można budować, przyjmując za oś odniesienia:
\begin{itemize}
    \item prąd -- dla połączenia szeregowego,
    \item napięcie -- dla połączenia równoległego.
\end{itemize}

\subsection*{Rozwiązanie}

\textbf{a) Połączenie szeregowe RC}

W połączeniu szeregowym przez oba elementy płynie ten sam prąd.
Przyjmujemy go jako wskaz odniesienia.

Napięcie rezystora jest zgodne w fazie z prądem, natomiast napięcie kondensatora
jest opóźnione względem prądu o $90^\circ$:

\[
\U_R=RI,
\qquad
\U_C=-\mathrm{j}\frac{I}{\omega C},
\qquad
\U=\U_R+\U_C.
\]

Moduł napięcia całkowitego:

\[
|\U|=I\sqrt{R^2+\left(\frac{1}{\omega C}\right)^2}.
\]

Kąt przesunięcia fazowego:

\[
\varphi=-\arctan\left(\frac{1}{\omega RC}\right).
\]

\begin{center}
\begin{tikzpicture}[>=stealth,scale=1.0]

\draw[->,gray] (-0.25,0)--(4.3,0) node[right] {$\Re$};
\draw[->,gray] (0,-2.1)--(0,1.0) node[above] {$\Im$};

\draw[->,thick,blue] (0,0)--(3.2,0)
node[midway,above] {$\U_R$};

\draw[->,thick,red] (3.2,0)--(3.2,-1.65)
node[midway,right] {$\U_C$};

\draw[->,very thick] (0,0)--(3.2,-1.65)
node[midway,below right] {$\U$};

\draw[->,thick] (0,0)--(1.35,0)
node[below] {$\I$};

\end{tikzpicture}
\end{center}

\textbf{b) Połączenie równoległe RC}

W połączeniu równoległym napięcie na obu elementach jest takie samo.
Przyjmujemy je jako wskaz odniesienia.

Prąd rezystora jest z nim zgodny w fazie, a prąd kondensatora wyprzedza
je o $90^\circ$:

\[
\I_R=\frac{U}{R},
\qquad
\I_C=\mathrm{j}\omega CU,
\qquad
\I=\I_R+\I_C.
\]

Moduł prądu całkowitego:

\[
|\I|=U\sqrt{\frac{1}{R^2}+(\omega C)^2}.
\]

Kąt przesunięcia fazowego:

\[
\varphi_I=\arctan(\omega RC).
\]

\begin{center}
\begin{tikzpicture}[>=stealth,scale=1.0]

\draw[->,gray] (-0.25,0)--(4.3,0) node[right] {$\Re$};
\draw[->,gray] (0,-0.35)--(0,2.8) node[above] {$\Im$};

\draw[->,thick,blue] (0,0)--(3.1,0)
node[midway,below] {$\I_R$};

\draw[->,thick,red] (3.1,0)--(3.1,1.9)
node[midway,right] {$\I_C$};

\draw[->,very thick] (0,0)--(3.1,1.9)
node[midway,above left] {$\I$};

\draw[->,thick] (0,0)--(1.35,0)
node[below] {$\U$};

\end{tikzpicture}
\end{center}

\textbf{Wniosek:}

W szeregowym RC napięcie całkowite jest sumą wskazów napięć,
a w równoległym RC prąd całkowity jest sumą wskazów prądów.

\newpage

% ==================================================
\section{Moduł impedancji szeregowego obwodu RC}
% ==================================================

\subsection*{Treść zadania}

Wyznaczyć moduł impedancji obwodu szeregowego RC dla częstotliwości
$f=1\,\mathrm{Hz}$ oraz $f=1\,\mathrm{MHz}$, przy:

\[
R=1\,\mathrm{k\Omega},
\qquad
C=1\,\mathrm{\mu F}.
\]

\subsection*{Najważniejsze wzory i teoria}

Dla połączenia szeregowego impedancje się sumują:

\[
\Z=R+\frac{1}{\jw C}
=R-\mathrm{j}X_C.
\]

Reaktancja kondensatora:

\[
X_C=\frac{1}{\omega C}
=\frac{1}{2\pi fC}.
\]

Moduł impedancji:

\[
|\Z|=\sqrt{R^2+X_C^2}.
\]

\subsection*{Dane i szukane}

\[
R=1000\,\Omega,
\qquad
C=10^{-6}\,\mathrm{F},
\]

\[
f_1=1\,\mathrm{Hz},
\qquad
f_2=10^6\,\mathrm{Hz}.
\]

Szukane:

\[
|\Z(f_1)|,\qquad |\Z(f_2)|.
\]

\subsection*{Rozwiązanie krok po kroku}

\textbf{1. Dla $f_1=1\,\mathrm{Hz}$.}

\[
\omega_1=2\pi f_1=2\pi\,\mathrm{rad/s}.
\]

\[
X_{C1}=
\frac{1}{2\pi\cdot1\cdot10^{-6}}
\approx159154.94\,\Omega.
\]

Podstawiamy do wzoru na moduł impedancji:

\[
|\Z_1|=
\sqrt{1000^2+159154.94^2}.
\]

\[
|\Z_1|\approx159158.08\,\Omega.
\]

\textbf{2. Dla $f_2=1\,\mathrm{MHz}$.}

\[
X_{C2}=
\frac{1}{2\pi\cdot10^6\cdot10^{-6}}.
\]

\[
X_{C2}\approx0.159155\,\Omega.
\]

\[
|\Z_2|=
\sqrt{1000^2+0.159155^2}.
\]

\[
|\Z_2|\approx1000.000013\,\Omega.
\]

\subsection*{Odpowiedź}

\[
\boxed{|\Z(1\,\mathrm{Hz})|\approx159.158\,\mathrm{k\Omega}}
\]

\[
\boxed{|\Z(1\,\mathrm{MHz})|\approx1.000000013\,\mathrm{k\Omega}}
\]

Przy małej częstotliwości dominuje reaktancja kondensatora,
a przy bardzo dużej częstotliwości jego reaktancja jest bardzo mała
w porównaniu z rezystancją.

\newpage

% ==================================================
\section{Moduł impedancji obwodu z rezystorem i cewką}
% ==================================================

\subsection*{Treść zadania}

Obliczyć moduł impedancji obwodu przedstawionego na schemacie
dla częstotliwości $f=1\,\mathrm{Hz}$ oraz $f=1\,\mathrm{MHz}$.

Wartości elementów:

\[
R_1=R_2=10\,\Omega,
\qquad
L=10\,\mathrm{mH}.
\]

\begin{center}
\begin{circuitikz}

\draw (0,2) to[short,o-] (0.7,2)
      to[R,l={$R_1=10\,\Omega$}] (2.6,2)
      -- (3.2,2);

\draw (3.2,2) to[L,l={$L=10\,\mathrm{mH}$}] (3.2,0);

\draw (3.2,2) -- (4.2,2)
      to[R,l_={$R_2=10\,\Omega$}] (4.2,0);

\draw (4.2,0)--(0,0)
      to[short,o-] (0,0);

\end{circuitikz}
\end{center}

\subsection*{Najważniejsze wzory i teoria}

Impedancja cewki:

\[
\Z_L=\jw L.
\]

Dla równoległego połączenia rezystora i cewki:

\[
\Z_{R_2\parallel L}
=
\frac{R_2\Z_L}{R_2+\Z_L}.
\]

Cały obwód składa się z $R_1$ połączonego szeregowo z tą gałęzią:

\[
\Z=R_1+\frac{R_2\jw L}{R_2+\jw L}.
\]

\subsection*{Rozwiązanie krok po kroku}

\textbf{1. Dla $f=1\,\mathrm{Hz}$.}

\[
\omega=2\pi\,\mathrm{rad/s}.
\]

\[
\Z_L=
\mathrm{j}2\pi\cdot1\cdot0.01.
\]

\[
\Z_L=\mathrm{j}0.062832\,\Omega.
\]

Impedancja gałęzi równoległej:

\[
\Z_{R_2\parallel L}
=
\frac{10(\mathrm{j}0.062832)}
{10+\mathrm{j}0.062832}.
\]

\[
\Z_{R_2\parallel L}
\approx
0.0003948+\mathrm{j}0.0628294\,\Omega.
\]

Całkowita impedancja:

\[
\Z_1
\approx
10.0003948+\mathrm{j}0.0628294\,\Omega.
\]

\[
|\Z_1|\approx10.000592\,\Omega.
\]

\textbf{2. Dla $f=1\,\mathrm{MHz}$.}

\[
\Z_L=
\mathrm{j}2\pi\cdot10^6\cdot0.01.
\]

\[
\Z_L=\mathrm{j}62831.853\,\Omega.
\]

\[
\Z_{R_2\parallel L}
=
\frac{10(\mathrm{j}62831.853)}
{10+\mathrm{j}62831.853}.
\]

\[
\Z_{R_2\parallel L}
\approx
9.99999975+\mathrm{j}0.00159155\,\Omega.
\]

Całkowita impedancja:

\[
\Z_2
\approx
19.99999975+\mathrm{j}0.00159155\,\Omega.
\]

\[
|\Z_2|\approx19.99999981\,\Omega.
\]

\subsection*{Odpowiedź}

\[
\boxed{|\Z(1\,\mathrm{Hz})|\approx10.000592\,\Omega}
\]

\[
\boxed{|\Z(1\,\mathrm{MHz})|\approx20.000000\,\Omega}
\]

Dla małej częstotliwości cewka ma małą impedancję.
Dla dużej częstotliwości impedancja cewki jest duża,
więc gałąź zachowuje się prawie jak sam rezystor $R_2$.

\newpage

% ==================================================
\section{Impedancja obwodu RLC i warunek rezonansu}
% ==================================================

\subsection*{Treść zadania}

Wyznaczyć impedancję obwodu, w którym cewka $L$ jest połączona
szeregowo z równoległym połączeniem rezystora $R$ i kondensatora $C$.

Określić warunek, jaki muszą spełniać wartości $R$, $L$ i $C$,
aby wystąpił rezonans, czyli aby obwód miał charakter czysto rezystancyjny.

\begin{center}
\begin{circuitikz}

\draw (0,2) to[short,o-] (0.6,2)
      to[L,l=$L$] (2.2,2)
      -- (3.2,2);

\draw (3.2,2) to[R,l=$R$] (3.2,0);

\draw (3.2,2)--(4.5,2)
      to[C,l=$C$] (4.5,0);

\draw (3.2,0)--(4.5,0)--(0,0)
      to[short,o-] (0,0);

\end{circuitikz}
\end{center}

\subsection*{Najważniejsze wzory i teoria}

Impedancja równoległego połączenia $R$ i $C$:

\[
\Z_{R\parallel C}
=
\left(
\frac{1}{R}+\jw C
\right)^{-1}.
\]

Po uproszczeniu:

\[
\Z_{R\parallel C}
=
\frac{R}{1+\jw RC}.
\]

Cewka jest połączona szeregowo z tą gałęzią, zatem:

\[
\Z(\omega)
=
\jw L+\frac{R}{1+\jw RC}.
\]

\subsection*{Rozwiązanie krok po kroku}

Usuwamy niewymierność z mianownika drugiego składnika:

\[
\Z(\omega)
=
\jw L+
\frac{R(1-\jw RC)}
{1+(\omega RC)^2}.
\]

Po rozdzieleniu na część rzeczywistą i urojoną:

\[
\boxed{
\Z(\omega)=
\frac{R}{1+(\omega RC)^2}
+
\mathrm{j}
\left(
\omega L-
\frac{\omega R^2C}
{1+(\omega RC)^2}
\right)
}
\]

Rezonans zachodzi, gdy część urojona impedancji jest równa zeru:

\[
\omega L-
\frac{\omega R^2C}
{1+(\omega RC)^2}=0.
\]

Dla $\omega\neq0$:

\[
L=
\frac{R^2C}
{1+(\omega_0RC)^2}.
\]

Stąd:

\[
(\omega_0RC)^2=
\frac{R^2C}{L}-1.
\]

Ostatecznie:

\[
\boxed{
\omega_0=
\sqrt{
\frac{1}{LC}
-
\frac{1}{R^2C^2}
}
}
\]

Częstotliwość rezonansowa:

\[
f_0=\frac{\omega_0}{2\pi}.
\]

Aby $\omega_0$ była rzeczywista i dodatnia, musi być spełniony warunek:

\[
\boxed{R^2C>L}
\]

czyli:

\[
\boxed{R>\sqrt{\frac{L}{C}}}
\]

W rezonansie impedancja jest czysto rzeczywista:

\[
\boxed{
\Z(\omega_0)=
\frac{R}{1+(\omega_0RC)^2}
=
\frac{L}{RC}
}
\]

\subsection*{Odpowiedź}

Warunek wystąpienia rezonansu:

\[
\boxed{R^2C>L}
\]

Częstotliwość rezonansowa:

\[
\boxed{
f_0=
\frac{1}{2\pi}
\sqrt{
\frac{1}{LC}
-
\frac{1}{R^2C^2}
}
}
\]

\newpage

% ==================================================
\section{Twierdzenie Thevenina dla dzielnika RC}
% ==================================================

\subsection*{Treść zadania}

Zastosować twierdzenie Thevenina do dzielnika RC.
Określić, jaki warunek musi być spełniony, aby napięcie Thevenina
$U_{AB}$ nie zależało od częstotliwości.

\begin{center}
\begin{circuitikz}

\draw (0,0) to[sV,l=$V_1$] (0,3)
      -- (1,3);

\draw (1.6,3) to[R,l=$R_1$] (1.6,1.5);
\draw (1.6,3)--(2.8,3);

\draw (2.8,3) to[C,l=$C_1$] (2.8,1.5);

\draw (1.6,1.5)--(3.8,1.5);

\draw (1.6,1.5) to[R,l=$R_2$] (1.6,0);
\draw (2.8,1.5) to[C,l=$C_2$] (2.8,0);

\draw (1.6,0)--(0,0);
\draw (2.8,0)--(0,0);

\draw (3.8,1.5) to[short,-o] (4.3,1.5)
node[right] {$A$};

\draw (3.8,0) to[short,-o] (4.3,0)
node[right] {$B$};

\end{circuitikz}
\end{center}

\subsection*{Najważniejsze wzory i teoria}

Każdy równoległy układ RC zastępujemy impedancją:

\[
\Z_1=
R_1\parallel\frac{1}{\jw C_1}
=
\frac{R_1}{1+\jw R_1C_1}.
\]

\[
\Z_2=
R_2\parallel\frac{1}{\jw C_2}
=
\frac{R_2}{1+\jw R_2C_2}.
\]

Napięcie Thevenina jest napięciem jałowym na zaciskach $A$--$B$.

Impedancję Thevenina wyznacza się po wyzerowaniu idealnego źródła napięcia.

\subsection*{Rozwiązanie krok po kroku}

Ponieważ $A$--$B$ jest wyjściem nieobciążonym,
napięcie jest napięciem dzielnika:

\[
\boxed{
\U_{\mathrm{Th}}=
\U_{AB}
=
\U_1\frac{\Z_2}{\Z_1+\Z_2}
}
\]

Po podstawieniu impedancji:

\[
\U_{AB}
=
\U_1
\frac{R_2(1+\jw R_1C_1)}
{R_1(1+\jw R_2C_2)+R_2(1+\jw R_1C_1)}.
\]

Po uporządkowaniu:

\[
\U_{AB}
=
\U_1
\frac{R_2+\jw R_1R_2C_1}
{R_1+R_2+\jw R_1R_2(C_1+C_2)}.
\]

Aby iloraz nie zależał od $\omega$,
współczynniki przy $\jw$ w liczniku i mianowniku muszą
być proporcjonalne do odpowiednich wyrazów stałych.

Otrzymujemy:

\[
\frac{R_1R_2C_1}{R_2}
=
\frac{R_1R_2(C_1+C_2)}{R_1+R_2}.
\]

Po uproszczeniu:

\[
\boxed{R_1C_1=R_2C_2}
\]

Jest to warunek równości stałych czasowych obu gałęzi.

Wtedy:

\[
\boxed{
\U_{AB}=
\U_1\frac{R_2}{R_1+R_2}
}
\]

czyli napięcie Thevenina nie zależy od częstotliwości.

Po zwarciu idealnego źródła napięcia obie impedancje
są połączone równolegle, dlatego:

\[
\boxed{
\Z_{\mathrm{Th}}=
\Z_1\parallel\Z_2
=
\frac{\Z_1\Z_2}{\Z_1+\Z_2}
}
\]

\subsection*{Odpowiedź}

Warunek niezależności napięcia Thevenina od częstotliwości:

\[
\boxed{R_1C_1=R_2C_2}
\]

Wtedy:

\[
\boxed{
\U_{AB}=
\U_1\frac{R_2}{R_1+R_2}
}
\]

\newpage

% ==================================================
\section{Wzmocnienie dwóch stopni wzmacniacza}
% ==================================================

\subsection*{Treść zadania}

Wzmocnienie dwóch stopni wzmacniacza napięciowego wynosi odpowiednio
$20\,\mathrm{dB}$ oraz $30\,\mathrm{dB}$.

Obliczyć całkowite wzmocnienie, stosując skalę decybelową
oraz bezwymiarową (V/V).

\subsection*{Najważniejsze wzory i teoria}

Dla wzmocnienia napięciowego:

\[
G_{\mathrm{dB}}=20\log_{10} A_v.
\]

Odwrotna zależność:

\[
A_v=10^{G_{\mathrm{dB}}/20}.
\]

Wzmocnienia kolejnych stopni połączonych kaskadowo mnożą się
w skali liniowej, natomiast w decybelach dodają się:

\[
A_{v,\mathrm{całk}}=A_{v1}A_{v2}.
\]

\[
G_{\mathrm{dB,całk}}
=
G_{\mathrm{dB1}}+G_{\mathrm{dB2}}.
\]

\subsection*{Rozwiązanie krok po kroku}

Całkowite wzmocnienie w decybelach:

\[
G_{\mathrm{dB,całk}}=20+30=50\,\mathrm{dB}.
\]

Wzmocnienia poszczególnych stopni w skali liniowej:

\[
A_{v1}=10^{20/20}=10.
\]

\[
A_{v2}=10^{30/20}=10^{1.5}\approx31.6228.
\]

Iloczyn:

\[
A_{v,\mathrm{całk}}
=
10\cdot31.6228
\approx316.228.
\]

Można to również obliczyć bezpośrednio:

\[
A_{v,\mathrm{całk}}
=
10^{50/20}
\approx316.228.
\]

\subsection*{Odpowiedź}

\[
\boxed{G_{\mathrm{dB,całk}}=50\,\mathrm{dB}}
\]

\[
\boxed{A_{v,\mathrm{całk}}\approx316.23\,\mathrm{V/V}}
\]

\newpage

% ==================================================
\section{Filtr górnoprzepustowy RC}
% ==================================================

\subsection*{Treść zadania}

Wyznaczyć zespolony współczynnik przenoszenia filtru
górnoprzepustowego RC i wyznaczyć częstotliwość graniczną.

\begin{center}
\begin{circuitikz}

\draw (0,1.5) to[short,o-] (0.7,1.5)
      to[C,l=$C$] (2.4,1.5)
      -- (3.2,1.5)
      to[R,l=$R$] (3.2,0)
      -- (0,0)
      to[short,o-] (0,0);

\draw (3.2,1.5) to[short,-o] (4,1.5)
node[right] {$U_{\mathrm{wy}}$};

\draw (0,1.5) node[left] {$U_{\mathrm{we}}$};

\end{circuitikz}
\end{center}

\subsection*{Najważniejsze wzory i teoria}

Współczynnik przenoszenia definiujemy jako:

\[
\Htr(\jw)=
\frac{\U_{\mathrm{wy}}}{\U_{\mathrm{we}}}.
\]

Dla dzielnika napięcia:

\[
\Htr(\jw)=
\frac{\Z_R}{\Z_C+\Z_R}.
\]

Impedancje:

\[
\Z_R=R,
\qquad
\Z_C=\frac{1}{\jw C}.
\]

\subsection*{Rozwiązanie krok po kroku}

Podstawiamy impedancje:

\[
\Htr(\jw)=
\frac{R}{R+\frac{1}{\jw C}}.
\]

Po uproszczeniu:

\[
\boxed{
\Htr(\jw)=
\frac{\jw RC}{1+\jw RC}
}
\]

Częstotliwość graniczna występuje dla:

\[
\omega_gRC=1.
\]

Stąd:

\[
\boxed{\omega_g=\frac{1}{RC}}
\]

oraz:

\[
\boxed{f_g=\frac{1}{2\pi RC}}
\]

Przy $f\ll f_g$ kondensator ma dużą reaktancję
i sygnał wyjściowy jest silnie tłumiony.

Przy $f\gg f_g$ reaktancja kondensatora jest mała
i $\Htr\to1$.

Jest to filtr górnoprzepustowy.

\subsection*{Odpowiedź}

\[
\boxed{
\Htr(\jw)=
\frac{\jw RC}{1+\jw RC}
}
\]

\[
\boxed{f_g=\frac{1}{2\pi RC}}
\]

\newpage

% ==================================================
\section{Współczynnik przenoszenia układu RC}
% ==================================================

\subsection*{Treść zadania}

Wyznaczyć współczynnik przenoszenia układu przedstawionego
na schemacie oraz omówić przypadki skrajne.

W układzie rezystor $3R$ jest połączony szeregowo z gałęzią,
w której rezystor $R$ i kondensator $C$ są połączone szeregowo.
Napięcie wyjściowe jest mierzone na zaciskach tej gałęzi.

\begin{center}
\begin{circuitikz}

\draw (0,1.7) to[short,o-] (0.7,1.7)
      to[R,l=$3R$] (2.8,1.7)
      -- (3.5,1.7)
      to[R,l=$R$] (3.5,0.75)
      to[C,l=$C$] (3.5,0)
      -- (0,0)
      to[short,o-] (0,0);

\draw (2.8,1.7) to[short,-o] (4.2,1.7);

\draw (3.5,0) -- (4.2,0)
      to[short,-o] (4.2,0);

\draw (0,1.7) node[left] {$U_{\mathrm{we}}$};

\draw (4.2,1.7) node[right] {$U_{\mathrm{wy}}$};

\end{circuitikz}
\end{center}

\subsection*{Najważniejsze wzory i teoria}

Układ jest dzielnikiem napięcia.

Impedancja gałęzi wyjściowej:

\[
\Z_{\mathrm{wy}}
=
R+\frac{1}{\jw C}.
\]

Współczynnik przenoszenia:

\[
\Htr(\jw)=
\frac{\Z_{\mathrm{wy}}}
{3R+\Z_{\mathrm{wy}}}.
\]

\subsection*{Rozwiązanie krok po kroku}

Podstawiamy impedancję gałęzi:

\[
\Htr(\jw)=
\frac{R+\frac{1}{\jw C}}
{3R+R+\frac{1}{\jw C}}.
\]

\[
\Htr(\jw)=
\frac{R+\frac{1}{\jw C}}
{4R+\frac{1}{\jw C}}.
\]

Mnożąc licznik i mianownik przez $\jw C$, otrzymujemy:

\[
\boxed{
\Htr(\jw)=
\frac{1+\jw RC}{1+\jw\,4RC}
}
\]

\textbf{Przypadek 1: częstotliwość dąży do zera.}

Dla $\omega\to0$:

\[
\lim_{\omega\to0}\Htr(\jw)=1.
\]

Kondensator zachowuje się jak przerwa w obwodzie.
Prąd jest bliski zeru, więc spadek napięcia na rezystorze $3R$
jest bliski zeru, a napięcie wyjściowe jest prawie równe wejściowemu.

\textbf{Przypadek 2: częstotliwość dąży do nieskończoności.}

Dla $\omega\to\infty$:

\[
\lim_{\omega\to\infty}\Htr(\jw)
=
\frac{R}{4R}
=
\frac14.
\]

Kondensator zachowuje się jak zwarcie,
więc pozostaje dzielnik rezystancyjny $3R$ i $R$.

Wysokoczęstotliwościowe wzmocnienie napięciowe wynosi $1/4$.

\subsection*{Wniosek}

Jest to układ o charakterystyce dolnoprzepustowej
z niezerowym wzmocnieniem dla wysokich częstotliwości.

Transmitancja zmienia się od $1$ do $1/4$.
Nie jest to idealny filtr dolnoprzepustowy tłumiący sygnał do zera.

\newpage

% ==================================================
\section{Filtr RL i jego odpowiednik RC}
% ==================================================

\subsection*{Treść zadania}

Wyznaczyć współczynnik przenoszenia przedstawionego filtru RL,
określić jego rodzaj oraz zastąpić go filtrem $R'C$,
wyznaczając odpowiednie wartości $R'$ i $C'$.

\begin{center}
\begin{circuitikz}

\draw (0,1.5) to[short,o-] (0.7,1.5)
      to[L,l=$L$] (2.6,1.5)
      -- (3.3,1.5)
      to[R,l=$R$] (3.3,0)
      -- (0,0)
      to[short,o-] (0,0);

\draw (3.3,1.5) to[short,-o] (4.1,1.5);

\draw (3.3,0) to[short,-o] (4.1,0);

\draw (0,1.5) node[left] {$U_{\mathrm{we}}$};

\draw (4.1,1.5) node[right] {$U_{\mathrm{wy}}$};

\end{circuitikz}
\end{center}

\subsection*{Najważniejsze wzory i teoria}

Napięcie wyjściowe jest mierzone na rezystorze $R$.

Dla dzielnika impedancyjnego:

\[
\Htr(\jw)=
\frac{\Z_R}{\Z_L+\Z_R}.
\]

Impedancje:

\[
\Z_R=R,
\qquad
\Z_L=\jw L.
\]

\subsection*{Rozwiązanie krok po kroku}

Podstawiamy impedancje:

\[
\Htr(\jw)=
\frac{R}{\jw L+R}.
\]

Po uproszczeniu:

\[
\boxed{
\Htr(\jw)=
\frac{1}{1+\jw\frac{L}{R}}
}
\]

Dla $\omega\to0$ współczynnik dąży do $1$,
natomiast dla $\omega\to\infty$ dąży do zera.

Jest to zatem \textbf{filtr dolnoprzepustowy RL}.

Dla filtru RC dolnoprzepustowego, w którym napięcie wyjściowe
jest mierzone na kondensatorze, współczynnik przenoszenia ma postać:

\[
\Htr_{RC}(\jw)=
\frac{1}{1+\jw R'C'}.
\]

Porównując oba wyrażenia, otrzymujemy warunek równoważności:

\[
\boxed{R'C'=\frac{L}{R}}
\]

Jest to jeden warunek na dwie niewiadome,
dlatego par $(R',C')$ jest nieskończenie wiele.

Można wybrać na przykład:

\[
R'=R.
\]

Wtedy:

\[
\boxed{C'=\frac{L}{R^2}}
\]

Można też wybrać dowolną dodatnią wartość $C'$
i obliczyć:

\[
R'=\frac{L}{RC'}.
\]

Oba filtry mają wtedy taki sam współczynnik przenoszenia.

\subsection*{Odpowiedź}

Współczynnik przenoszenia:

\[
\boxed{
\Htr(\jw)=
\frac{1}{1+\jw\frac{L}{R}}
}
\]

Rodzaj filtru: dolnoprzepustowy RL.

Warunek równoważności filtru RC:

\[
\boxed{R'C'=\frac{L}{R}}
\]

Przykładowy dobór elementów:

\[
\boxed{R'=R,\qquad C'=\frac{L}{R^2}}
\]

% ==================================================
\section*{Podsumowanie}
% ==================================================

\begin{itemize}

    \item Szeregowy obwód RC:
    impedancja $R+1/(\jw C)$.

    \item Równoległy obwód RC:
    prądy gałęzi sumują się wskazowo.

    \item Dla rezonansu w obwodzie z cewką szeregową
    i równoległym RC wymagane jest $R^2C>L$.

    \item Napięcie dzielnika z równoległymi gałęziami RC
    nie zależy od częstotliwości, gdy
    $R_1C_1=R_2C_2$.

    \item Wzmocnienia kaskadowe dodają się w dB,
    a mnożą w skali liniowej.

\end{itemize}

\end{document}

